\newpage

\section{Zaključak}

Realizacija ovog projekta pokazala je da je moguće uspješno integrirati moderan infotainment sustav u starije vozilo koristeći kombinaciju vlastitog hardverskog modula i open-source softverskih alata. Ključni izazovi bili su razumijevanje i implementacija VAN protokola, osiguranje stabilnog napajanja te izrada pouzdanog firmwarea za ESP32 mikroupravljač. Pažljivim dizajnom tiskane pločice, korištenjem provjerenih komponenti i pisanjem modularnog koda, postignuta je robusna i stabilna komunikacija s vozilom. \\

Prednosti rješenja su otvorenost i mogućnost proširenja funkcionalnosti, niska cijena u odnosu na komercijalne alternative te potpuna kontrola nad softverom i hardverom. Sustav se pokazao pouzdanim tijekom testiranja – prikaz telemetrije, statusa vozila i reprodukcija multimedije radili su bez većih odstupanja od tvorničkih pokazivača. Najveći izazovi, poput dekodiranja VAN okvira i implementacije režima štednje energije, uspješno su savladani korištenjem ESP-IDF okvira i pažljivim upravljanjem događajima i zadacima.\\

Ipak, postoje područja za daljnje unaprjeđenje. Potrebno je dodatno optimizirati potrošnju energije u stanju mirovanja kako bi se maksimalno smanjilo pražnjenje akumulatora te doraditi korisničko sučelje na Raspberry Pi računalu radi intuitivnijeg korištenja. Budući razvoj mogao bi uključivati potpunu integraciju s postojećim zaslonom u vozilu, dodavanje bežične povezivosti (Wi-Fi/Bluetooth sinkronizacija s pametnim telefonom) te naprednije funkcije poput OTA (Over-the-Air) nadogradnje softvera. Time bi se sustav dodatno približio funkcionalnosti i fleksibilnosti modernih tvorničkih infotainment rješenja.

